\documentclass[10pt]{article}
\usepackage[a4paper]{geometry}
\usepackage{amsmath}

% a macro, to show that \newcommand works
\newcommand{\knit}{\textsf{knitting}}
\newcommand{\term}[1]{\emph{#1}}

\title{Typesetting on a Worker Pool}
\author{knit-tex \and the \knit{} runtime}
\date{\today}

\begin{document}
\maketitle

\begin{abstract}
This document was set by knit-tex, a small TeX-flavoured typesetter written in
plain JavaScript. Every paragraph is broken into lines with the Knuth--Plass
algorithm and hyphenated with Liang's patterns, the same two ideas that sit at
the heart of \TeX. Each paragraph is an independent task, so the line breaking
runs on a \knit{} worker pool while the main thread only parses the source,
stacks lines onto pages and writes the file.
\end{abstract}

\section{Why paragraphs are a good unit of work}
\label{sec:why}

A paragraph is the natural grain of a typesetting job. Within it, every
decision depends on every other: moving one word to the next line changes how
loose or tight all the following lines can be, which is why a good line
breaker looks at the paragraph as a whole instead of filling one line at a
time. Across paragraphs, nothing is shared at all. That shape---tightly
coupled inside, completely independent outside---is exactly what a worker pool
wants.\footnote{Pages are the second unit of work: once lines are stacked onto
pages, every page's content stream is generated and compressed in parallel too.}

The cost per task is small. A typical paragraph of prose takes somewhere
between a tenth of a millisecond and a few milliseconds to set, depending on
its length and on how many hyphenation passes it needs. Work that fine-grained
is where the overhead of a worker call matters most: if dispatching a task
costs as much as running it, there is nothing to gain. That is the property of
\knit{} this example leans on.

\subsection{The algorithm in one breath}

Knuth and Plass model a paragraph as a sequence of \term{boxes} (things that
are printed), \term{glue} (space that can stretch or shrink) and
\term{penalties} (places where a break is allowed but has a cost). For a line
from breakpoint $a$ to breakpoint $b$ with natural width $w$, total stretch
$Y$ and total shrink $Z$, the adjustment ratio is
\begin{equation}
  r = \frac{L - w}{Y} \quad\text{when } w < L, \qquad r = \frac{L - w}{Z} \quad\text{when } w > L,
  \label{eq:ratio}
\end{equation}
where $L$ is the line width. The badness of the line is $100\,|r|^3$, and the
demerits of a break combine badness with the penalty $p$ of the breakpoint:
\[
  d = (\ell + b)^2 + p^2 \qquad (p \geq 0).
\]
The optimal set of breaks minimises the sum of demerits over the whole
paragraph, which the algorithm finds with dynamic programming over a small set
of \term{active nodes}. Lines whose ratio in~\eqref{eq:ratio} falls below
$-1$ would have to shrink more than the glue allows, so they are never
chosen unless nothing else is possible.

\subsection{What \TeX\ adds on top}

The paper describes the core; \TeX\ adds the details that make the output look
right, and knit-tex follows them closely:
\begin{itemize}
  \item three passes: first without hyphenation at a strict tolerance, then
        with hyphenation, then an emergency pass that always succeeds;
  \item four fitness classes, so a very loose line is not placed next to a
        very tight one;
  \item extra demerits for two hyphenated lines in a row, and for a
        hyphen on the second-to-last line of a paragraph;
  \item \TeX's own rule for hopeless paragraphs: when the last active node is
        about to disappear, keep it and accept an overfull line, so a single
        unbreakable word can never stop a document from being typeset.
\end{itemize}

\section{The pipeline}

\begin{enumerate}
  \item \textbf{Parse} the source on the main thread into blocks: paragraphs,
        headings, lists, verbatim text and displayed equations.
  \item \textbf{Set} every paragraph on the pool. A task receives styled runs
        of text and returns positioned segments, one array per line.
  \item \textbf{Stack} lines onto pages on the main thread, keeping headings
        with what follows them and avoiding widows and orphans.
  \item \textbf{Render} each page's content stream on the pool, compressed
        with deflate.
  \item \textbf{Write} the PDF: fonts, pages, bookmarks and the cross-reference
        table.
\end{enumerate}

The tasks themselves are tiny wrappers, and they live in their own module so
that workers only load what they run:
\begin{verbatim}
import { task } from "knitting";
import { setParagraph } from "./paragraph.mjs";
import { renderPageStream } from "./render.mjs";

export const setParagraphTask = task({
  f: (spec) => setParagraph(spec),
});
export const renderPageTask = task({
  f: (page) => renderPageStream(page),
});
\end{verbatim}

\subsection{Fonts without font files}

PDF viewers must provide fourteen standard fonts, so knit-tex embeds none: it
measures text with the Adobe font metrics for Times, Courier and Symbol,
including kerning pairs such as ``AV'', ``To'' and ``We''. Characters outside
the Windows-1252 encoding fall back to the Symbol font when it has them, which
is how Greek letters such as $\alpha$, $\beta$ and $\Omega$, and operators such
as $\leq$, $\infty$ and $\nabla$, reach the page. Accented names come through
too: Erdős, Gödel, Dvořák, Ångström, naïve café.

\subsection{Lists, quotes and descriptions}

\begin{description}
  \item[Boxes] are words, or the syllables of a word that can be hyphenated.
  \item[Glue] is the space between words. Its stretch and shrink are a half
        and a third of the space width, as in Computer Modern.
  \item[Penalties] mark hyphenation points (cost 50) and forced breaks.
\end{description}

The benchmark novel opens with a sentence that every line breaker has to get
right:
\begin{quote}
It is a truth universally acknowledged, that a single man in possession of a
good fortune, must be in want of a wife.\\
\hspace*{\fill}---Jane Austen, \emph{Pride and Prejudice} (1813)
\end{quote}

Nested lists use the same label styles as \LaTeX:
\begin{enumerate}
  \item First level, numbered.
  \begin{enumerate}
    \item Second level, lettered.
    \item Another one, with an em dash---like this---and an en dash for
          ranges such as pages 12--14.
  \end{enumerate}
  \item Back to the first level.
\end{enumerate}

\section{Results}

Section~\ref{sec:why} argued that paragraphs are fine-grained work. The
benchmark in the repository typesets \emph{Pride and Prejudice} (about 120{,}000
words, two thousand paragraphs) on the main thread and on pools of one to three
workers, and checks that every configuration produces a byte-identical
PDF.\footnote{Determinism is a design goal: the writer stamps no dates, and
floating-point layout runs the same code in every thread.} The output you are
reading was produced the same way.

\begin{thebibliography}{9}
\bibitem{kp} Donald E. Knuth and Michael F. Plass. Breaking paragraphs into
  lines. \emph{Software: Practice and Experience}, 11(11):1119--1184, 1981.
\bibitem{liang} Franklin M. Liang. \emph{Word Hy-phen-a-tion by Com-put-er}.
  PhD thesis, Stanford University, 1983.
\end{thebibliography}

\end{document}
